\chapter{データセットと前処理}
\section{対象データと分割}
UAV3DCropのCorn、Oat、Wheat、Soybeanから、RGBと対応深度GTの512組、計1024ファイルを使用した。記録された転送量は7,850,330,776 byteである。データは第1回実験時にGoogle Driveへ保存され、第2・3回で再利用された。本報告書作成時にデータは再取得・再監査していない。
\begin{table}[H]\centering\caption{実験で使用した分割}\label{tab:split}
\begin{tabular}{llll}\toprule 用途&撮影日&組数&使用\\\midrule
学習&Day019, Day028, Day040&384&学習\\
検証&Day047&128&各runのvalidation\\
独立test&未取得&319&未使用\\
過去評価&過去の52組&52&モデル選択に未使用\\\bottomrule\end{tabular}
\end{table}
撮影日単位で分割し、同じ日の類似画像がtrain/validation間へ混在するリスクを抑えた。ただし、未知日・未知圃場への汎化を保証するものではない。独立testはモデル選択完了まで未使用に保留された。
\section{GTと前処理}
RGBと深度は元解像度3956$\times$5280 pixelで管理された。train GT最大値は62.15625 m、99.9 percentileは52.90625 mである。train最大値を包含する教師・評価上限として$G=65$ mを設定した。これはvalidation/testによる選択ではない。無効値は0で、GTが有限かつ正かつ65 m以下の画素だけを有効とした。
学習は518$\times$518 cropを入力とする。検証はアスペクト比を保持して短辺518、長辺約686 pixelへresizeし、予測を元サイズへ復元して評価した。元短辺3956 pixelから518への倍率は約7.64であり、元画像14 pixel間隔はモデル入力上約1.83 pixelとなる。ViTの14 pixel patchを目安とすると元画像上約107 pixelに対応する。DPTは多解像度特徴を融合するものの、縮小で葉・茎・株境界の情報が弱まる可能性がある。
\section{深度と草丈の違い}
GT Zはカメラ光軸方向に沿う距離である。Obliqueでは地面平面も画像内で奥行き変化をもつ。カメラ俯角のある鉛直植物は、見かけのZ変化が高さと単純比例しない。草丈評価には共通世界座標で地面と頂部を対応づけ、鉛直成分を計算する必要がある。したがって20 cmのZ差は垂直草丈20 cmを意味しない。